\subsection{Encoder layers: attention, FFN, residual, LayerNorm}
\label{sec:expl-encoder}

\begin{intuitbox}
One encoder layer does two jobs: (1)~let every token look at every other token (self-attention);
(2)~nonlinearly transform each token row alone (feed-forward network).
Residuals keep old information; LayerNorm keeps numbers stable across depth.
\end{intuitbox}

\begin{whybox}
The slide shows the attention formula. Examiners may ask what else sits inside each of the twelve layers---this subsection answers that.
\end{whybox}

BioBERT applies the same block \(L=12\) times. The update from layer \(\ell-1\) to layer \(\ell\) is
\begin{align}
  Z^{(\ell)} &= \mathrm{LayerNorm}\!\bigl(H^{(\ell-1)} + \mathrm{MultiHeadAttention}(H^{(\ell-1)})\bigr),
  \label{eq:Zell}\\
  H^{(\ell)} &= \mathrm{LayerNorm}\!\bigl(Z^{(\ell)} + \mathrm{FFN}(Z^{(\ell)})\bigr),
  \label{eq:Hell}\\
  \mathrm{FFN}(x) &= W_2\,\mathrm{ReLU}(W_1 x + b_1) + b_2.
  \label{eq:FFN}
\end{align}
In BERT-base / BioBERT-base:
\(W_1 \in \mathbb{R}^{768 \times 3072}\), \(W_2 \in \mathbb{R}^{3072 \times 768}\),
\(b_1 \in \mathbb{R}^{3072}\), \(b_2 \in \mathbb{R}^{768}\).
The inner width \(3072 = 4\times 768\) gives the FFN capacity to transform each row nonlinearly.

\begin{center}
\small
\begin{tabular}{@{}lp{9cm}@{}}
\toprule
Component & What it does \\
\midrule
Multi-head attention & Twelve parallel heads; each learns different token relationships \\
FFN & Expand to 3072 dims with ReLU, then project back to 768 \\
LayerNorm & Rescale activations for stable deep training \\
Residual (\(+\)) & Add sub-layer input to output so earlier signal is not erased \\
\bottomrule
\end{tabular}
\end{center}

\paragraph{Stacking twelve times.}
\[
  H^{(0)} \xrightarrow{\text{Layer 1}} H^{(1)} \xrightarrow{\text{Layer 2}} \cdots \xrightarrow{\text{Layer 12}} H^{(12)} = H^{(L)}.
\]
Rough intuition (not a hard theorem): shallow layers emphasise words and short phrases;
middle layers combine symptoms (``headache'' + ``feverish'' + ``weak'');
deep layers encode broader urgency patterns. All \(M\) rows are processed in parallel at each layer.

\begin{remark}
The \texttt{[CLS]} row of \(H^{(12)}\) is exactly \(h_{\mathrm{[CLS]}}\) from Equation~\eqref{eq:hcls}.
\end{remark}

\begin{saybox}
``Each layer: attend, add \& normalise, feed-forward, add \& normalise. Do that twelve times. Then read \texttt{[CLS]}.''
\end{saybox}

%----------------------------------------------------------------------
\subsection{Self-attention in detail (slide 18)}
\label{sec:expl-attn}

\begin{intuitbox}
Attention lets every token decide how much to listen to every other token.
``severe'' should strengthen ``headache''; ``no'' should weaken ``pain''.
\end{intuitbox}

\begin{whybox}
This is the heart of transformers. Know \(Q,K,V\), the scale \(\sqrt{d_k}\), and what \(\alpha_j\) means.
\end{whybox}

\paragraph{Four computational steps.}
From the current hidden matrix \(H\):
\begin{enumerate}
  \item \textbf{Project} into queries, keys, and values.
  \item \textbf{Score} pairwise compatibility.
  \item \textbf{Normalise} scores into attention weights.
  \item \textbf{Combine} value rows into context-aware outputs.
\end{enumerate}

\begin{align}
  Q &= H W_Q, \qquad K = H W_K, \qquad V = H W_V,
  \label{eq:QKV}\\
  S &= \frac{QK^{\top}}{\sqrt{d_k}},
  \label{eq:S}\\
  A &= \mathrm{softmax}(S),
  \label{eq:A}\\
  \mathrm{output} &= A V.
  \label{eq:AV}
\end{align}
Equivalently, the compact formula on the slide \citep{vaswani2017attention}:
\begin{equation}
  \mathrm{Attention}(Q,K,V)
  = \mathrm{softmax}\!\left(\frac{QK^{\top}}{\sqrt{d_k}}\right)V,
  \qquad d_k = 64.
  \label{eq:attn}
\end{equation}

\paragraph{What \(Q\), \(K\), \(V\) mean in words.}
\begin{itemize}
  \item \(Q\) (query): ``what am I looking for at this position?''
  \item \(K\) (key): ``what do I offer as a matchable label?''
  \item \(V\) (value): ``what content should be mixed if I match?''
\end{itemize}

\paragraph{Why divide by \(\sqrt{d_k}\)?}
Dot products of \(d_k\)-dimensional vectors grow with \(d_k\).
Without scaling, softmax saturates (weights become nearly \(0\) or \(1\)), gradients vanish, and learning stalls.
The factor \(\sqrt{d_k}\) keeps scores in a healthy range (\(d_k=64\) per head in BioBERT base).

\paragraph{Attention weights \(\alpha_j\).}
For one query row, the softmax over keys assigns weight \(\alpha_j\) to token \(j\):
\begin{equation}
  \alpha_j = \frac{\exp(e_j)}{\sum_{k=1}^{m}\exp(e_k)}, \qquad \sum_{j=1}^{m}\alpha_j = 1,
  \label{eq:alpha}
\end{equation}
where \(e_j\) is the scaled score for key \(j\) and \(m\) is the sequence length (masked pads do not participate).
\textbf{Meaning:} \(\alpha_j\) is how much this position listens to token \(j\).
Symptom tokens typically receive higher weight than fillers such as ``I'' or ``and''.

\paragraph{Multi-head attention.}
Twelve heads run in parallel with different \(W_Q,W_K,W_V\).
Outputs are concatenated and projected back to \(768\) dimensions so heads can specialise
(e.g.\ symptom co-occurrence vs negation vs intensity words).

\paragraph{Toy attention pattern (intuition only).}
When the query is \texttt{[CLS]}, illustrative (not measured) weights might look like:
\begin{center}
\small
\begin{tabular}{@{}lc@{}}
\toprule
Token & Relative weight \\
\midrule
head / \texttt{\#\#ache} & high \\
fever / \texttt{\#\#ish} & high \\
weak & medium--high \\
I, have, a, and & low \\
\texttt{<PAD>} & zero (masked) \\
\bottomrule
\end{tabular}
\end{center}
So the summary vector emphasises clinical wording, not filler words.

\begin{saybox}
``\(QK^\top\) scores who listens to whom; divide by \(\sqrt{64}\); softmax makes weights that sum to one; multiply by \(V\). Twelve heads, then stack twelve layers.''
\end{saybox}
